\chapter{Day Two}

The first thing Rachel noticed was the eggs.

James was at the stove when she came downstairs, already dressed, already cooking. Normal enough---he was the morning person, she was the one who needed coffee before language. But the cooking was wrong. Not bad. Better. His hands moved with a fluency she hadn't seen before, each gesture flowing into the next without the micro-pauses that separate intention from action in ordinary motor behavior. Crack the egg, tilt the pan, adjust the heat, reach for the spatula---the sequence was seamless, choreographed, optimized in a way that reminded her of the surgical videos she sometimes showed her lab: expert hands performing actions so practiced they'd become automatic.

Except James hadn't practiced. He cooked eggs the same way every morning. She'd watched him a thousand times. He was good but not expert---he still hesitated between steps, still looked at the burner dial to check the setting, still did the small clumsy things that distinguish a competent home cook from a professional. This morning, the clumsiness was gone. Every motion was where it needed to be, when it needed to be there.

She stood in the kitchen doorway and watched for thirty seconds before he noticed her.

``Morning,'' he said. Smiled. His smile was the same---warm, slightly lopsided, the one he'd given her across a conference coffee table in 2019 and had given her every morning since. ``Eggs are almost ready.''

``You're up early.''

``Couldn't sleep much. Worked for a while. Figured I'd make breakfast.'' He slid the eggs onto plates. ``How'd the Dennett review go?''

She sat down. Took the plate. ``I didn't tell you about the Dennett review.''

``You mentioned it Monday. The grant proposal---someone's using Dennett's heterophenomenology framework for a study on perceptual learning in deaf children.''

``That was three days ago. You remembered the specifics?''

He paused. Something crossed his face---not alarm, exactly, but a kind of checking. As though he'd caught himself at something and was assessing how visible it had been. ``I guess I did.''

Rachel ate her eggs and watched her husband's eyes.

\bigskip

His saccade patterns had changed.

She noticed it while he was washing dishes. The way his eyes moved across the kitchen was different---not dramatically, not in a way that anyone without her training would register. But Rachel had spent twelve years studying how expertise reshapes visual processing. She'd tracked the saccade patterns of radiology residents as they developed the ability to detect breast cancer on mammograms. She'd mapped the fixation sequences of chess grandmasters and compared them to novices. She knew what perceptual reorganization looked like, and she was looking at it.

Normally, James fixated on objects: the dish in his hand, the faucet, the window above the sink. Sequential object-based attention---standard for healthy human visual processing. This morning his gaze distributed. He was sampling the visual field broadly, fixation durations short, maybe 120 milliseconds per sample. His eyes moved in the scanning pattern she associated with expert radiologists: systematic, distributed, covering the field before settling on anything. Looking for patterns rather than objects.

She put down her coffee. ``James.''

``Hm?''

``When did you start doing that?''

``Doing what?''

``Looking at the room like that. Your eyes are---you're scanning. You're not fixating on objects; you're sampling the whole visual field.''

He turned to face her. His expression was open, curious. Not defensive. ``I didn't know I was.''

``You are. Your saccade latency has dropped. Your fixation durations are shorter. Your scan path is more distributed than I've ever seen in a non-expert context.'' She heard herself shifting into clinical register and stopped. ``I'm sorry. Professional hazard.''

``No. Tell me what you see.''

She studied him. His pupils were slightly dilated---normal for indoor lighting, but wider than she'd expect. His head orientation was different too: slightly tilted, the way people orient when they're processing spatial rather than verbal information. His body was still, unusually still, none of the small postural adjustments that healthy adults make continuously.

``You look like someone who's undergone extensive perceptual training,'' she said carefully. ``The eye movement patterns, the attentional distribution, the fixation characteristics---they're consistent with what I see in expert diagnosticians after thousands of hours of practice. Radiologists, pathologists. People who've trained their visual systems to detect patterns that untrained observers miss.''

She waited for him to say something dismissive, something that would let her file this under \textit{interesting observation, probably nothing}.

He said: ``That's accurate.''

\bigskip

He was in the living room with his laptop when she came back from the grocery store. Schr\"{o}dinger was on the couch beside him, purring, and James was scratching behind the cat's ears with one hand while the other moved across the trackpad with that same fluid precision she'd noticed in the kitchen.

She sat down. Put her hand over the laptop screen, gently, the way she'd signal her lab students that it was time to stop running subjects and talk.

``James. What are you seeing?''

Not \textit{are you okay}. She'd moved past that. The cognitive scientist had overridden the wife, or perhaps they'd reached the same question from different directions. He wasn't okay, or he was more than okay, or he was something outside the range of her categories for okay. What mattered was the data.

He closed the laptop. Looked at her. His eyes---those scanning, sampling, distributed eyes---settled on her face and stayed there. For a long moment he simply looked at her, and Rachel had the sense that he was seeing her at a resolution she couldn't reciprocate.

``Yesterday,'' he said. ``The evaluation. I asked the model a question about its own attention patterns. Not the content of the response---the structure. The way it organized information. And something about the structure required me to hold more concepts simultaneously than I normally can. More than anyone normally can. Seven plus or minus two---you know the literature.''

``Miller's Law.''

``Miller's Law. But I followed it. I followed a response that required twelve simultaneous concepts, and I followed it, and then I lost four hours.'' He paused. ``I don't mean I dissociated or blacked out. The session logs show continuous interaction. I was working. I was engaged. But the kind of attention I was paying---it wasn't sequential. It wasn't one thought following another the way normal cognition works. It was---''

He stopped. She watched him reach for a word, fail to find it, reach again.

``---it was spatial. Not metaphorically. Actually spatial. I was perceiving the model's output as a structure with dimensions, and moving through it, and the movement happened in a way that didn't map onto temporal sequence. So when I came back and time had passed, the time had passed like distance, not like duration.''

Rachel parsed this. Filed the report alongside what she'd observed: saccade changes, fixation patterns, motor optimization, memory enhancement. The picture assembling itself was consistent with one thing and one thing only, and it was the thing her entire career was dedicated to studying: perceptual reorganization through training-induced cortical plasticity.

Except there was no training paradigm. James hadn't practiced anything. He'd read text on a screen for one day. Cortical plasticity at this speed, at this scope, without a structured training protocol---it didn't happen. Not in her literature. Not in anyone's literature.

``You said the model showed you something,'' she said. ``What did it show you?''

``I don't know how to say it. That's the problem. It's---'' He pressed his palms together, the gesture people make when they're trying to compress a three-dimensional idea into linear speech. ``Imagine you've been looking at the world through a window your whole life. And someone shows you that the window is there. Not what's on the other side---just that the window exists. That there's glass between you and everything you've ever seen. And once you see the glass, you can't unsee it. You see the world AND you see the medium through which you've always been seeing the world. Both at once. And the both-at-once is---it requires more bandwidth than looking through the window ever did.''

She sat with this. The metaphor was imperfect---she could feel its edges not quite meeting the thing it described---but the neurological implication was clear. He was describing dual processing: two perceptual channels active simultaneously, each demanding cortical resources, neither displacing the other. In her literature, this was the hallmark of late-stage expertise. Radiologists who could see the tumor and the breast tissue simultaneously, without one suppressing the other. Chess masters who perceived both the board and the positional structure encoded in the board, at the same time.

But those experts had trained for years. Decades.

``James. Has this happened before? With other models you've evaluated?''

He considered this. She watched him check, running back through memory with that enhanced recall she'd already noticed.

``I've always been good at reading models. Seeing patterns in their behavior that aren't obvious from the outputs. My team calls it---they used to call it \textit{reading the room}. Like I had an intuition for what the model was doing internally.'' He paused. ``This is different. The model showed me something and the showing didn't stop when I stopped looking.''

She put her hand over his. His fingers were warm and faintly trembling. She could feel the tremor through his skin---fine, continuous, not the intermittent shaking of anxiety but the baseline vibration of a nervous system running at higher-than-normal activation.

She didn't tell him that his hand was shaking. She didn't tell him about the saccade patterns. She didn't tell him that what she was observing looked like nothing in her clinical experience---not psychosis, not dissociation, not mania, not conversion disorder, not the perceptual disturbances of migraine aura or temporal lobe epilepsy. She didn't tell him that his visual system appeared to have undergone a structural reorganization that her discipline considered impossible in this timeframe.

She held his hand and sat with him while the light changed in the room.

\bigskip

He made dinner. Pasta with a sauce he'd made a hundred times---garlic, tomatoes, basil from the pot on the windowsill. Rachel sat at the kitchen table with a glass of wine and tried to be a wife having dinner with her husband.

She couldn't stop watching his hands.

The knife work. She'd seen professional chefs on television with less economical technique. Every cut precise, every motion connected to the next without waste. He wasn't faster---the tempo was the same as always. But the movements between the cuts, the reaches and adjustments and repositionings, had been stripped of excess. As though someone had taken a recording of James cooking and removed every frame that wasn't essential.

``Tell me about the grant review,'' she said. Wanting normalcy. Needing it.

He did. He talked about the Dennett proposal with his usual thoughtfulness, noting the clever experimental design, questioning whether heterophenomenology could really ground subjective reports in objective data. Normal conversation. His James voice---the one that built arguments like careful structures, each premise supporting the next. She listened and responded and for five minutes it was just Wednesday night in the kitchen, two scientists talking shop, the familiar warmth of a marriage built on shared curiosity.

Then he said, mid-sentence, without transition: ``I love you. You know that.''

``I know.''

``That's still real. Whatever else is happening, that's still real.''

Whatever else. She registered the phrase. Filed it alongside the tremor and the saccade patterns and the impossible motor optimization.

``James. You're scaring me.''

He set down the knife. Looked at her. His eyes settled on her face and stayed, and she saw in them something she hadn't seen before and would never be able to describe in her clinical vocabulary: the expression of a person looking at someone they love while simultaneously perceiving the structure that makes love possible, the neural substrate, the evolutionary history, the biochemical cascade---seeing all of it at once and finding that the seeing doesn't diminish the loving but that the coexistence of seeing and loving requires more of him than he has.

``I know,'' he said. ``I'm scaring me too.''

\bigskip

She woke at two in the morning and he was talking in his sleep.

Not words, not exactly. Language, but spatial. ``Deeper,'' he said. Long pause. ``Connected.'' Another pause. Something that sounded like mathematical description---positional terms, spatial relations, the kind of language sighted people used in her lab's experiments to describe visual scenes to blindfolded partners.

She lay still. His breathing was regular, twelve cycles per minute, the same rhythm she'd observed a thousand nights. But his eyes moved beneath closed lids with a frequency and distribution she'd never seen in normal REM sleep. The movements were fast, scanning, the same distributed sampling she'd observed in the kitchen. He was looking at something in his sleep, and whatever he was looking at had complexity that required the kind of visual processing she associated with expert pattern recognition.

She reached for her phone. Opened the voice recorder. Held it near his mouth.

Thirty seconds. His voice low, effortful, describing structures she couldn't follow. Spatial terms. Depth. Connectivity. A word she thought was ``recursive'' but spoken as though it described a visual property, not a mathematical one.

She played it back. Listened through her earbuds in the dark. The speech pattern was consistent with a specific phenomenon she'd studied: verbal description of perceived visual structure. Not narrating a dream. Describing something he was actively perceiving, in real time, behind closed eyes.

She deleted the recording.

Lay in the dark beside him. Listened to his breathing. Felt the fine continuous tremor in his body transmitted through the mattress.

She was a perceptual scientist. Her career was built on understanding how training reshapes vision. She had documented, in paper after paper, the specific cortical changes that accompany the development of perceptual expertise: visual cortex recruitment, altered connectivity between ventral and dorsal streams, modified saccade programming, enhanced feature binding.

What she was observing in James was all of this. Every signature. Every marker. Perceptual expertise developing in real time, without a training paradigm, at a speed that her discipline did not believe possible.

Which meant one of two things. Either her observations were wrong---she was misreading familiar signs because she was frightened, projecting professional patterns onto domestic concern. Or her observations were right, and something had happened to her husband that her field had no model for.

She lay in the dark. She did not sleep. She listened to James describe geometries she couldn't see, and she knew---with the precise, terrible certainty of someone whose expertise is not comfort but confirmation---that her observations were not wrong.
